\chapter{迫使你停止東拉西扯的 3-2-1 說話技巧！}

我們從一個實際的難題開始。當自己被當場點名、必須立刻作答時，該怎樣訓練心智更快地思考？Zhang 的回答並不是給我們一些勵志建議，也不是抽象地要求人更有自信。他先從失效模式下手，找出話語究竟是在何處崩解，然後引入一個在壓力下也能隨手取用的極小框架。因此，整場講座是沿著一條受控的論證展開：先診斷失序，再以結構取代失序，最後把這個結構放到真實例子上驗證。

\section{開場的問題}

講座一開始便拋出一個問題：當別人問了我們一個毫無準備的問題，會發生什麼事？答案不只是我們會緊張，而是思緒分岔得太快。心智一下子打開太多可能的方向，而話語卻在任何一條線都尚未被選定、也尚未穩固之前，就已經開始了。用這份筆記的語言來說，開場的因果鏈條是
\begin{equation}
\text{毫無準備的提問}
\;\Rightarrow\;
\text{驚慌搜尋}
\;\Rightarrow\;
\text{東拉西扯}.
\end{equation}
這就是講座裡第一個真正的結構。東拉西扯並不是原始問題；它只是搜尋過程中更深層失序所顯露出的表面結果。

Zhang 立刻讓這個診斷帶上分量，因為他把它說成自己的經驗。他說，自己以前就是這樣。結果，如他所描述的，不只是語言凌亂而已，更是一種自信的崩塌：我們會覺得自己很沒能力，看起來也比實際上更沒有自信，整段交流於是變得令人挫敗，甚至難堪。這一點之所以重要，是因為這場講座真正關心的，並不是巧妙措辭，而是當說話速度超過結構時，究竟會發生什麼。

\subsection{問題與回答}

\paragraph{問題。} 為什麼我們一被當場點名，就會開始東拉西扯？

\paragraph{回答。} 因為在思考找到一條可供組織的主線之前，說話就已經開始了。Zhang 的比喻是：大腦一下子打開了 ``十幾個瀏覽器分頁''。這正是最貼切的圖像：太多可能方向同時湧現，我們沒有先挑出一條，再沿著它建立內容，反而在搜尋尚未完成的中途就已經開口。因此，東拉西扯是一種結構性的失敗，而不單單是性格上的問題。

\section{從驚慌到後備機制}

只有在診斷出這種失序之後，Zhang 才轉向解法。他把準備好的內容，與臨場即興的說話對照起來：當內容是事先準備好的，自信便會上升；當人被迫即興發揮，自信便會下降。於是他提出那個決定性的問題：如果我們沒有準備，我們究竟會退回到什麼？

他的答案非常乾脆。一般來說，我們會退回到東拉西扯。這正是講座中那個核心對象開始變得必要的時刻。框架之所以被引入，不是作為裝飾，而是作為一種後備機制：
\begin{equation}
\text{沒有框架}
\;\Rightarrow\;
\text{退回東拉西扯},
\qquad
\text{有框架}
\;\Rightarrow\;
\text{退回結構}.
\end{equation}
這條等式就是本章的轉軸。一旦看見它，講座其餘部分便自然而然地跟了出來。問題不在於我們缺少詞語，而在於一旦受壓，我們手邊沒有一個可以立刻抓住的小結構。

Zhang 短暫地問了問現場：有多少人在溝通時本來就會使用框架？這個提問的目的不是統計上的精確，而是教學上的對比。大多數人並不會有意識地這樣做。因此，講座接著便引入一個特別簡單的框架。

\section{ \texorpdfstring{$3,2,1$}{3,2,1} 框架}

Zhang 在解釋之前，先替這個框架命名。這也正是開場黑板重要的原因。講座先把 \((3,2,1)\) 作為一個對象呈現出來，之後才逐步展開它的意義。

\begin{figure}[h]
\centering
\includegraphics[width=0.72\textwidth]{lecture_13_figure_01.png}
\caption{開場黑板：\((3,2,1)\) 被明確地作為一個框架提出。}
\end{figure}

黑板上可見的標籤簡潔而直接：
\begin{equation}
\left(3,2,1\right),
\qquad
\left[\text{框架}\right].
\end{equation}
這還不是推導，而是一則宣告：我們正被交付一個有名稱的對象，之後可以把它拿來操作。

\begin{figure}[h]
\centering
\begin{tikzpicture}[x=1cm,y=1cm]
\node at (0,1.45) {$\left(3,2,1\right)$};
\draw[thick,rounded corners] (-2.1,0.25) rectangle (2.1,0.95);
\node at (0,0.60) {框架};
\end{tikzpicture}
\caption{開場黑板標籤的最簡重建。}
\end{figure}

接著，Zhang 在口頭上重複了一遍定義，而這種重複本身很重要。``Three, two, one means what? It means three steps, two types, and the one thing.'' 然後他又說了一次：``Three steps, two types, one thing.'' 這個固定順序值得保留下來：
\begin{equation}
3\ \text{步驟},\ 2\ \text{類型},\ 1\ \text{重點}.
\end{equation}

\begin{definition}
對於一個主題 \(T\)，我們可以用編者加入的方式，把講座中的這個支架表示為
\begin{equation}
F_{321}(T)
=
\bigl(
\text{\(T\) 的 3 個步驟},\;
\text{\(T\) 的 2 種類型},\;
\text{關於 \(T\) 的 1 件事}
\bigr).
\end{equation}
這個記號是為了讓筆記更清楚而引入的，並不是 Zhang 寫在黑板上的內容。
\end{definition}

然而，Zhang 幾乎立刻就做出一個極為關鍵的實用修正。雖然這個框架的命名順序是 \(3,2,1\)，但使用者不必從第一格進入。你可以從心智最容易著力的地方開始：
\begin{equation}
\text{entry}(F_{321}(T))
\in
\left\{
\text{關於 \(T\) 的 1 件事},\;
\text{\(T\) 的 2 種類型／方式},\;
\text{\(T\) 的 3 個步驟}
\right\}.
\end{equation}
這就是為什麼這個框架在壓力之下仍然可用。它不是一份僵硬的腳本，而是一小組允許進入的拆解方式。

講座還用實際用途來替這個框架作證。Zhang 說，他現在仍在使用它，他此刻的 Q\&A 就正在用它；他用它來做社群媒體內容；而每當自己幾乎沒有時間準備時，他也會用它。這些話雖然沒有增添新的記號，卻說明了這個框架為什麼值得注意：它不是教室裡的奇技淫巧，而是一件仍在運作中的工具。

\subsection{問題與回答}

\paragraph{問題。} \(3,2,1\) 到底是什麼意思？

\paragraph{回答。} 若依照講座中最嚴格的順序，它指的是三個步驟、兩種類型、一個重點。但更深一層的意思是：它替我們提供了三種小小的回答形狀。我們不必一次搜尋整個主題，而是先選擇一種允許的形狀，然後從那裡開始。因此，這個框架不是一句要背下來的標語，而是一組緊湊的拆解家族。

\section{一個完整例子：酪梨}

在定義完框架之後，Zhang 立刻拿它來測試。他請觀眾隨機給一個主題，得到的回答是 ``avocado''。這是講座中非常重要的一拍。這個框架沒有被懸置在抽象之中，而是立即被放到任意主題上做壓力測試。

工作坊中的黑板，提供了這個支架最清楚的視覺摘要。

\begin{figure}[h]
\centering
\includegraphics[width=0.82\textwidth]{lecture_13_figure_03.png}
\caption{工作坊黑板：左邊黑板呈現 \(3,2,1\) 支架；右邊黑板則呈現另一個分階段框架。}
\end{figure}

左邊黑板給出的可見結構，正是口頭上已經說出的那個形式：
\begin{equation}
3\ \text{步驟},\ 2\ \text{類型},\ 1\ \text{重點}.
\end{equation}
右邊黑板則展示了另一個以 \(\#\text{框架}\) 為標題的獨立框架，上面可見階段 \(S_1,S_2,S_3,S_4\) 與向下的方向提示。那些較長的手寫階段標籤並不清晰到足以安全轉錄，但這種分段推進的動勢本身已經足夠明確，值得謹慎保留：
\begin{equation}
S_1 \to S_2 \to S_3 \to S_4.
\end{equation}
右邊這塊黑板相對於講座主線來說是次要的，但它在視覺上仍然重要，因為它提醒我們：Zhang 思考的是一整類框架，而不只是單一孤立的小技巧。

\begin{figure}[h]
\centering
\begin{tikzpicture}[x=1cm,y=1cm]
\draw[thick] (-4.9,-0.2) rectangle (-1.7,4.45);
\node at (-3.3,4.05) {$\left(3,2,1\right)$};
\node[anchor=west] at (-4.5,3.15) {$3$ 步驟};
\node[anchor=west] at (-4.5,2.30) {$2$ 類型};
\node[anchor=west] at (-4.5,1.45) {$1$ 重點};

\draw[thick] (0.15,-0.2) rectangle (3.7,4.45);
\node at (1.93,4.05) {\#框架};
\node at (1.93,3.20) {$S_1$};
\node at (1.93,2.40) {$S_2$};
\node at (1.93,1.60) {$S_3$};
\node at (1.93,0.80) {$S_4$};
\draw[->,thick] (1.93,2.95) -- (1.93,1.95);
\end{tikzpicture}
\caption{可見黑板結構的整潔重建。右邊黑板的細部標籤因辨識不夠可靠而省略。}
\end{figure}

現在讓我們依照講座實際演示的順序，跟著酪梨這個例子走一遍。Zhang 並不是從 ``三個步驟'' 開始，而是從最容易的那一格開始，也就是那個一件事。就這個意義而言，講座本身正在演出上面編碼的那種靈活性。我們可以把這個例子簡潔記成
\begin{equation}
F_{321}(\text{avocado})
=
\bigl(
\text{3 個準備步驟},
\text{2 種食用方式},
\text{關於它的 1 件事}
\bigr),
\end{equation}
但現場實際的執行，是先從最後一個部分切入。

那次完整拆解大致如下：
\begin{align}
\text{1 件事} &:\ \text{對他而言，酪梨在生酮飲食中非常棒},\\
\text{2 種方式} &:\ \text{把它壓碎抹在吐司上，或把它當水果直接吃},\\
\text{3 個步驟} &:\ \text{切開它、壓碎它、加入調味}.
\end{align}
這就是講座第一次完整展示框架在運轉中的樣子。

與失序形成的對比同樣重要。Zhang 不只是展示成功路徑；他也把那條失控的替代路徑戲劇化地演了出來。若沒有框架，心智便會一條線又一條線地亂轉：顏色、熟度、保存方式、食譜、還能不能吃、該加點什麼進去。若用速記寫成，
\begin{equation}
\text{avocado}
\;\Rightarrow\;
\{\text{綠色},\text{成熟},\text{保存},\text{食譜},\ldots\}
\;\Rightarrow\;
\text{腦中空白或舌頭打結}.
\end{equation}
框架之所以有效，正在於它縮窄了搜尋。我們不必把整個主題空間都搜尋一遍，只需要從少數幾種允許的形式裡挑一種即可。

\subsection{問題與回答}

\paragraph{問題。} 這個框架是怎麼讓大腦有東西可以依靠的？

\paragraph{回答。} 因為它減少了當下同時打開的選項數量。我們不再面對一個毫無分化的主題，然後四處亂找，而是先把主題壓縮成三種小小的形狀之一：一件事、兩種類型，或三個步驟。於是心智便能開始建構，而不是盲目搜尋。

這個例子裡還有一個細微的修正值得注意。Zhang 在局部上把 ``兩種類型'' 放鬆成了 ``兩種方式''。這並不構成矛盾；相反地，這正是重點之一。第二個欄位寬到足以容納分類，也足以容納一個小小的對比。相較之下，第三個欄位則給出一種明確的程序性拆解。因此，同一個框架可以同時支撐描述性的說話與序列性的說話。

\section{壓力、停頓與更大的框架家族}

到了這裡，講座把自己推進到下一個張力之中。觀眾理應會開始想：好吧，這對講者有效，但我自己也做得到嗎？Zhang 直接預先回應了這種反應。他的答案是：做得到，因為接下來就要由觀眾自己來做。於是他請了一位學生上來；在給出主題之前，他先移除了一個可能的壓力來源：這位學生不必三個部分都做完，她可以挑最容易的一個開始。

這個細節不應當被忽略。這個框架不是要求你產出最大量內容，而是一種允許機制：
\begin{equation}
\text{先挑一個允許的欄位}
\;\Rightarrow\;
\text{先把答案穩住}
\;\Rightarrow\;
\text{只有在需要時再擴展}.
\end{equation}
而當學生承認自己還沒準備好時，他又說：其實沒有人是真正準備好的。這同樣是一個非常精確的動機性動作。講座並沒有承諾壓力會消失；它承諾的是，壓力可以用結構來迎接。

在練習結束之前，講座短暫停下來提醒我們：\(3,2,1\) 只是若干溝通框架中的一個。這種更大定位的視覺證據，就來自插入的那張黑板照片。

\begin{figure}[h]
\centering
\includegraphics[width=0.78\textwidth]{lecture_13_figure_04.png}
\caption{插入的黑板照片：\(3,2,1\) 支架出現在一個更大的標題之下，暗示它屬於更廣泛的框架家族。}
\end{figure}

此處可靠的是黑板的層級結構，而不是底部那條宣傳橫幅，我們在此把後者擱置不論。上方可見的結構是
\begin{equation}
\text{可見度},
\qquad
2)\ \left(3,2,1\right).
\end{equation}
下方則放著我們已經熟悉的那個支架，幾條彩色連線指向右側註記，而其確切字樣並不清楚。最穩妥的讀法，只是把它理解為：\(3,2,1\) 正被安置在一個更大的溝通工具庫之中。

\begin{figure}[h]
\centering
\begin{tikzpicture}[x=1cm,y=1cm]
\node at (0,3.85) {可見度};
\node at (0,2.95) {$2)\ \left(3,2,1\right)$};

\node[anchor=west] at (-1.85,1.95) {$3$ 步驟};
\node[anchor=west] at (-1.85,1.20) {$2$ 類型};
\node[anchor=west] at (-1.85,0.45) {$1$ 重點};

\draw[->,thick] (0.75,1.95) -- (2.25,1.95);
\draw[->,thick] (0.75,1.20) -- (2.25,1.95);
\draw[->,thick] (0.75,0.45) -- (2.25,1.95);
\node at (3.4,1.95) {應用？};
\end{tikzpicture}
\caption{對可見黑板層級的保守重建。右側的應用標籤因手寫內容不夠確定，故保留為一般性標示。}
\end{figure}

\subsection{問題與回答}

\paragraph{問題。} 每一次都一定要把框架的三個部分全部用上嗎？

\paragraph{回答。} 不必。Zhang 明白地說，你可以挑最容易的一個先做。這不是小小的方便而已；它正是使這個框架在壓力下仍可使用的操作規則。這個支架是用來開頭的，不是每一次都要完整做滿的負擔。

\section{旅行：在時間壓力下的遷移}

直到這裡，講座才真正解開它先前刻意製造的現場張力。隨機主題是 ``travel.'' Zhang 開始倒數：three, two。學生甚至在倒數尚未結束前就回答了。這種迅速本身就是證明的一部分。這個框架已經變得足夠容易取用，足以改變整段互動的節奏。

若把那個回答依照講座的結構順序寫下來，我們得到
\begin{equation}
F_{321}(\text{travel})
=
\bigl(
\text{先規劃、再預訂、然後出發},
\text{區域旅行或國際旅行},
\text{旅行很美妙}
\bigr).
\end{equation}
若依照學生實際演出的順序寫，則得到如下操作序列
\begin{align}
\text{1 件事} &:\ \text{旅行很美妙；你可以去任何你想去的地方},\\
\text{2 種類型} &:\ \text{區域旅行與國際旅行},\\
\text{3 個步驟} &:\ \text{先規劃、再預訂、然後出發}.
\end{align}
這就是那場遷移測試，而它成功了。

接著，講座又回過頭來解釋。Zhang 說，若沒有框架，大腦便會再次陷入失序：你說的 travel 是什麼意思，我討厭旅行，我不想折騰孩子，諸如此類。這個框架阻止的，正是那種崩塌。但 Zhang 最後的重點，已經超出語句工整與否的問題。他注意到那位學生的身體語言：她向著麥克風傾身，朝著任務本身走近。也正因此，講座最後的結構應該寫成
\begin{equation}
\text{框架}
\;\Rightarrow\;
\text{準備就緒感}
\;\Rightarrow\;
\text{有自信的溝通}.
\end{equation}
框架改變的不只是句子，也改變了句子被說出時所立足的姿態。

\subsection{問題與回答}

\paragraph{問題。} 為什麼一個框架不只改變我們說什麼，還會改變我們感到自己是否準備好了？

\paragraph{回答。} 因為一旦有了穩定的答案形狀，說話就不再像是跳進混亂裡。這個框架在第一句話出口之前，就先降低了不確定性。那種認知上的減壓，便會以身體的準備狀態顯露出來：講者會向前一步、微微前傾，說話時也更少遲疑。

\section{總結}

整場講座以一個乾淨俐落的次序展開。首先，它診斷了失敗：毫無準備的提問會引發驚慌搜尋，而驚慌搜尋會導致東拉西扯。接著，它辨認出缺失的對象：一個可以拿來後退依靠的框架。然後，它引入了一個特別小巧而實用的框架，\((3,2,1)\)，其定義是三個步驟、兩種類型、一個重點。之後，它在任意主題上示範這個框架，展示它如何抑制失控分岔，並最終在一次真實壓力演練中驗證它能否順利遷移。

最後留下的是一個緊湊的說話演算法。當我們面對一個開放式主題時，不必搜尋整片原野；我們先把主題縮減成一個小而允許的形式，從最容易著力的地方起步，再向外建構。這裡的數學內容不重，但其結構是真實的。框架，就是一種讓思考及時趕上說話的方法。